\subsection*{刚体模型、位置与姿态}

真实固体受力后都会发生形变，因为构成物体的微观粒子之间距离会改变。刚体是假设任意两质点
距离在运动中保持不变的理想模型。只有当形变相对于物体尺寸很小，并且形变对所研究的运动没有
显著影响时，才可以采用这一模型。刚体并不意味着物体“非常硬”，而是说明本题忽略弹性波、
局部压缩和内部振动；碰撞、细长杆折断或高速旋转问题若涉及这些效应，就必须重新检查刚体近似。

质心的三个坐标只能确定刚体在空间中的位置，不能确定它的朝向。例如一本书的质心不动，仍可
绕质心任意转动。三维刚体还需要三个独立参数描述姿态，因此自由刚体共有六个自由度。Euler
角给出一种统一的姿态描述：先绕空间固定轴转动，再绕新产生的节线转动，最后绕刚体固连轴转动。
三次转动的先后次序不可交换，也不能把三个角的时间导数直接当作角速度在三根固定直角轴上的
分量。若题目只研究平面运动，姿态可以缩减为一个转角，角速度也只剩垂直于运动平面的分量。

按运动方式可把常见刚体运动分为三类。平动中刚体各点速度相同，任意固连线段始终平行于自身；
定轴转动中有一条固连直线保持不动，各点绕该轴作圆周运动；平面平行运动则同时包含质心平动和
绕质心的转动。车轮滚动、杆在平面内滑动以及圆盘与地面接触都属于第三类。刚体的复杂运动不是
新的运动学，而是质心平动与转动两种基本运动的组合。

\subsection*{速度分布与瞬时转动中心}

刚体上任一点 $P$ 的速度为
\[
  \vect v_P=\vect v_C+\vect\omega\times\vect r_{CP}.
\]
对同一刚体上的任意两点 $A,B$，相减得到
\[
  \vect v_B-\vect v_A=\vect\omega\times\vect{AB},
  \qquad
  (\vect v_B-\vect v_A)\cdot\vect{AB}=0.
\]
这说明两点相对速度必与两点连线垂直，是检验速度图是否正确的直接条件。在平面运动中，若能
找到速度为零的点 $P$，刚体在该瞬间的速度场等价于绕 $P$ 转动：
\[
  \vect v_Q=\vect\omega\times\vect{PQ}.
\]
因此每一个速度箭头都必须与 $PQ$ 垂直，长度与 $PQ$ 成正比。瞬时中心只是某一时刻速度场的
几何工具，它下一时刻通常会移动，也不表示那里存在固定铰链。固定地面上的纯滚动轮，其接触点
当时速度为零，所以可作为瞬时中心；轮缘上下一时刻与地面接触的是另一个质点。

加速度不能照搬瞬时中心的速度结论。即使接触点瞬时速度为零，它仍可能有非零法向加速度。
刚体两点的加速度关系为
\[
  \vect a_B-\vect a_A
  =\vect\beta\times\vect{AB}
  +\vect\omega\times(\vect\omega\times\vect{AB}),
\]
第二项是向心项。把瞬时中心误当作固定转轴并令它的加速度为零，会在滚动加速度问题中得到错误
结果。

\subsection*{动量方程与转动方程}

刚体总动量只取决于质心运动：
\[
  \vect P=M\vect v_C,
  \qquad
  \sum\vect F_{\rm ext}=M\vect a_C.
\]
同一外力无论作用在何处，对质心加速度的贡献相同；但它对转动状态的影响由作用线决定。关于
质心写角动量定理，平面刚体常用
\[
  \sum M_C=I_C\beta.
\]
若改为固定惯性轴，也可写相同形式并使用关于该轴的转动惯量。对任意加速点直接写
$\sum M=I\beta$ 一般不成立，必须补上由参考点运动产生的项。

圆形刚体沿倾角为 $\theta$ 的斜面向下纯滚动，是平动方程、转动方程和运动学约束同时使用的
典型例子。设半径为 $R$，关于质心的转动惯量为 $I$，沿斜面方向有
\[
  Mg\sin\theta-f=Ma_C,
  \qquad fR=I\beta,
  \qquad a_C=\beta R.
\]
解得
\[
  a_C=\frac{Mg\sin\theta}{M+I/R^2},
  \qquad
  f=\frac{I}{I+MR^2}Mg\sin\theta.
\]
不同刚体即使外半径与总质量相同，也会因质量分布不同而具有不同的 $I$，从而以不同加速度
下滚。原图把不同内部结构的圆形刚体并列画出，正是为了强调“外轮廓相同不等于滚动响应相同”。

纯滚动条件不能预先假定成立。求出所需静摩擦后还要检查 $|f|\le\mu_sN$；若不满足，应改用
滑动摩擦，并删除 $a_C=\beta R$。静摩擦方向也不是永远沿斜面向上，应根据若无摩擦时接触点
将向哪边相对滑动来判断。外加转矩、传送带或偏心质量都可能使摩擦方向反转。

\subsection*{刚体动能、功率与冲量}

刚体动能可以分解为
\[
  E_k=\frac12Mv_C^2+\frac12I_C\omega^2.
\]
第一项是质心平动动能，第二项是相对质心的转动动能。外力系统对刚体的瞬时功率也相应分解：
\[
  P=\vect F_{\rm ext}\cdot\vect v_C+\vect M_C\cdot\vect\omega.
\]
固定粗糙面上纯滚动时，接触点瞬时速度为零，理想静摩擦对整个刚体的总功率为零；但它仍会
改变平动与转动动能之间的分配。不能由“静摩擦不做总功”推出“静摩擦对运动没有作用”。

碰撞时间很短时，力随时间的积分形成冲量，力矩随时间的积分形成冲量矩：
\[
  \Delta\vect P=\int\vect F\,\dd t,
  \qquad
  \Delta\vect L_O=\int\vect M_O\,\dd t.
\]
选碰撞点为矩心常能消去未知接触冲量；选固定铰点为矩心可消去铰链冲量。打击中心问题要求既
产生所需转动，又使铰点冲量为零。对均匀细杆，打击中心位置由质量、转动惯量和铰点到质心距离
共同决定，并不是凭几何中心猜得。烟囱倒地和长杆碰撞还需要考察截面弯矩、剪力；刚体整体方程
给出运动后，再对截取部分列方程，才能判断哪里最先达到材料强度极限。

\subsection*{瞬心法能做什么、不能做什么}

平面运动刚体在某一瞬间总可以写成绕瞬时速度为零点 $P$ 的转动。若已知两点 $A,B$ 的速度方向，
分别过 $A,B$ 作速度的垂线，两条垂线的交点就是瞬心；随后有
\[
 v_A=\omega PA,\qquad v_B=\omega PB.
\]
这既给出速度大小之比，也给出每个速度必须垂直于从瞬心连向该点的半径。修图时若连线方向已经
确定，速度箭头就必须严格与该连线垂直，不能只凭“看起来顺着转动”估计。

瞬心是速度场的概念，不是固定铰点，也不保证该点加速度为零。刚体各点加速度关系为
\[
 \vect a_A=\vect a_B+\vect\alpha\times\vect r_{A/B}
 +\vect\omega\times(\vect\omega\times\vect r_{A/B}).
\]
其中切向项垂直于连线，法向项沿连线指向转动中心。因此不能把速度图整体旋转九十度就当作
加速度图。纯滚动接触点瞬时速度为零，却通常具有非零加速度；若接触点离开地面，这个结论还会
随约束改变。

\subsection*{滚动问题的完整求解顺序}

处理斜面滚动时，先规定质心平动正方向和角加速度正方向，再分别写
\[
 \sum F_\parallel=Ma_C,\qquad \sum M_C=I_C\alpha.
\]
只有题目说明无滑动或静摩擦足够大，才补上 $a_C=\alpha R$。联立以后求得的摩擦力带符号；若
结果为负，表示实际方向与最初假设相反。最后用 $|f|\le\mu_sN$ 检查约束是否自洽。若不满足，
改用动摩擦 $f_k=\mu_kN$，此时平动和转动方程仍成立，但 $a_C=\alpha R$ 必须撤去。

从静止沿同一斜面下滚、且都保持纯滚动的物体，质心加速度可写为
\[
 a_C=\frac{g\sin\theta}{1+\kappa},\qquad
 \kappa=\frac{I_C}{MR^2}.
\]
实心球、实心圆柱、薄环的 $\kappa$ 依次不同，故到达斜面底端的先后由质量分布决定，与物体总
质量无关。比较图中若只画相同外圆而不呈现内部质量分布，就会遗漏题目真正考查的变量。

\subsection*{整体运动与内部受力}

刚体模型规定质点间距离不变，却不意味着内部作用力为零。细杆、烟囱或转动构件在运动中必须靠
内部拉力、压力、剪力和弯矩维持几何约束。求断裂位置时，先用整体方程得到角速度和角加速度，
再在待考察截面处截开一部分，对该部分列平动、转动方程，求出截面内力随位置的分布。最大应力
位置未必等于最大外力位置，也未必在几何中点。

能量法擅长连接两个状态，却通常不能单独给出约束反力；动力学方程能给瞬时力和加速度，却可能
需要更多未知量。竞赛题中常先用能量求 $\omega$，再用 $\sum\vect F=M\vect a_C$ 和
$\sum\vect M_C=I_C\vect\alpha$ 求接触力或铰链力。两套方法各司其职，比强行只用一种方法更
稳妥，也更便于用量纲、极限和方向相互校验。

\subsection*{自由度与约束的计数}

空间中一个自由刚体有三个平动自由度和三个转动自由度。固定一点后，平动自由度被约束，只剩
三个转动自由度；固定轴转动只剩一个自由度；平面运动通常有两个平动自由度和一个转动自由度。
无滑动接触又在速度层面加入约束。列方程前先计数自由度和独立约束，可以预判需要多少广义坐标，
也能发现是否重复使用了同一个约束条件。

Euler 角只是描述空间姿态的一种坐标选择，并不等同于三个互相独立、始终固定方向的角速度分量。
有限转动一般不可交换，先绕一轴再绕另一轴与相反次序得到的姿态不同；只有无穷小转动在一阶近似
下可以按矢量相加。处理三维刚体时，应把姿态参数、瞬时角速度和角动量的方向区分开来，尤其对
非球对称刚体，$\vect L$ 未必与 $\vect\omega$ 平行。
